\documentclass[11pt]{article}

\usepackage[margin=1in]{geometry}
\usepackage{booktabs}
\usepackage{amsmath,amssymb}
\usepackage{hyperref}
\usepackage{xcolor}

\hypersetup{
  colorlinks=true,
  linkcolor=blue,
  citecolor=blue,
  urlcolor=blue
}
\setlength{\emergencystretch}{3em}

\title{LLM Alpha-Agent Replication Report}
\author{Sasha Cui}
\date{July 9, 2026}

\begin{document}
\maketitle

\begin{abstract}
This report evaluates whether recent LLM- and agent-based alpha-mining papers and
public repositories contain reproducible, economically meaningful equity return
signals once they are placed into a common monthly USA-equity asset-pricing
benchmark. I audit the source universe in two layers. First, I ask whether each
public repository can itself produce a dated return stream from approved local
inputs. Second, when direct code is absent or unusable, I translate the stated
economic idea into an explicit JKP-USA proxy and test it against FF5+Mom,
JKP132, and TextBenchmark/NewsFactor-style benchmarks. The conclusion is
negative for repo-level alpha discovery: direct public code does not produce a
valid FF5+Mom-beating JKP return stream, and the few in-spirit survivors are
overwhelmingly classic value, quality, profitability, momentum, and low-risk
composites.
\end{abstract}

\section{Question}

The empirical question is deliberately narrow. Let \(r_t\in\mathbb{R}\) be the
monthly excess return of a candidate strategy at month \(t\), and let
\(f_t\in\mathbb{R}^{K}\) be the benchmark factor return vector. A finance-agent
paper or repository is treated as economically convincing only if it yields a
dated return stream \(r_t\) that survives common benchmark tests:
\[
  r_t = \alpha + \beta^\top f_t + \epsilon_t.
\]
The relevant claim is not whether LLMs can help financial research. The relevant
claim is whether public finance-agent repositories have discovered new alpha in a
form that survives standard asset-pricing controls and portfolio-usefulness
tests.

\section{Benchmark Design}

The audit has two layers:
\begin{enumerate}
  \item Direct-code audit: can the public repository itself produce dated monthly
  USA-equity returns from approved local inputs?
  \item In-spirit JKP replication: if direct code is absent or unusable, can the
  stated economic idea be translated into an explicit JKP-USA proxy and tested
  against the same benchmark environment?
\end{enumerate}

Valid candidate and benchmark returns use only approved local research inputs
from the JKP/USA return panel and the existing TextBenchmark/NewsFactor factor
panel. China/A-share returns, yfinance or live downloads, paper-shipped return
streams, official French downloads, and external return CSVs are not counted as
valid metrics.

The first benchmark is FF5+Mom:
\[
  f_t =
  (\mathrm{MKT}_t,\mathrm{HML}_t,\mathrm{ME}_t,\mathrm{INV}_t,
  \mathrm{OP}_t,\mathrm{MOM}_t)^\top,
\]
implemented from JKP-built USA factors. The stricter additive benchmark evaluates
candidates against CAPM, JKP132, and the TextBenchmark/NewsFactor sleeve.

\section{Metrics}

For a candidate return \(r_t\), the reported diagnostics are annualized Sharpe,
annualized alpha, HAC/Newey-West alpha \(t\)-statistic, appraisal ratio, GRS
statistic and \(p\)-value, factor-span Sharpe lift, and long-only book Sharpe
improvement. The appraisal ratio is
\[
  AR=\sqrt{12}\frac{\widehat{\alpha}}{\widehat{\sigma}_{\epsilon}},
\]
and the factor-span Sharpe lift is
\[
  \Delta SR = \sqrt{SR_{\mathrm{old}}^2 + AR^2} - SR_{\mathrm{old}}.
\]
The strict FF5+Mom rule requires positive annualized alpha, positive appraisal
ratio, HAC alpha \(t>1.96\), positive span lift, and GRS rejection at the 5\%
level. The strict additive rule further requires survival against
CAPM+JKP132+TextBenchmark, positive long-only delta Sharpe, and at least a 1\%
optimized candidate weight.

\section{Headline Counts}

\begin{table}[h]
\centering
\begin{tabular}{lr}
\toprule
Quantity & Count \\
\midrule
Literature inventory & 42 paper rows \\
Code-link inventory & 31 rows \\
Unique source references & 55 \\
Direct cloned/repo audit & 14 repos \\
Code-link repo evidence pass & 28 GitHub repos \\
Direct repo-code rows with valid JKP-only numeric returns & 1 \\
Direct repo-code rows beating FF5+Mom & 0 \\
Source refs translated into JKP-USA proxies & 51 of 55 \\
JKP-USA candidate proxies backtested & 62 \\
Strict FF5+Mom beaters among in-spirit proxies & 8 of 62 \\
Strict additive candidates after TextBenchmark and JKP132 & 3 of 62 \\
\bottomrule
\end{tabular}
\caption{Headline replication counts. Full ledgers are tracked in
\texttt{paper\_runs/} and summarized in \texttt{report.md}.}
\end{table}

\section{Interpretation}

The direct-code conclusion is negative. QuantEvolver is the only repository with
valid JKP-only numeric seed-proxy metrics, and its selected declared-direction
seed has Sharpe \(0.316\), annualized alpha \(0.33\%\), HAC \(t=0.651\), appraisal
ratio \(0.149\), GRS \(p=0.4778\), and factor-span lift \(0.010\). It does not
beat FF5+Mom.

The in-spirit proxy layer finds eight strict FF5+Mom beaters, but these are not
credible evidence of new agentic alpha. They are almost entirely value, quality,
profitability, safety, momentum, low-leverage, and low-volatility composites.
After adding JKP132 and TextBenchmark/NewsFactor, only three of sixty-two
candidates remain strict positives. Those survivors are Greenblatt-style value
and profitability, a multi-guru value, quality, and distress ensemble, and a
value, profitability, momentum, and quality-minus-risk composite. The full
CAPM+JKP132+TextBenchmark span explains roughly 93.5\% to 95.8\% of their monthly
return variance.

The right interpretation is therefore not that LLMs are irrelevant. It is that
unverified finance-agent alpha claims should be treated as research-process
artifacts until they produce dated returns that survive common benchmarks. This
supports a more disciplined role for LLMs in financial research: constructing new,
structured, auditable information objects from economically rich source material,
then testing those objects in a standard portfolio framework.

\section{Reproducibility}

The canonical Markdown report is \texttt{report.md}. The compact tracked ledgers,
verdicts, source tables, and performance figures live under \texttt{paper\_runs/}.
The most important artifacts are:
\begin{itemize}
  \item \texttt{paper\_runs/REPOSITORY\_FF5MOM\_METRICS.md}
  \item \texttt{paper\_runs/idea\_replications/PAPER\_DERIVED\_REPLICATION\_LEDGER.md}
  \item \path{paper_runs/performance_analysis/textbenchmark/TEXTBENCHMARK_PERFORMANCE_REPORT.md}
  \item \texttt{paper\_runs/performance\_analysis/figures/FIGURE\_MANIFEST.md}
\end{itemize}

\begin{thebibliography}{99}
\bibitem{carhart1997} Carhart, Mark M. 1997. On Persistence in Mutual Fund Performance. \emph{Journal of Finance} 52(1), 57--82.
\bibitem{cochrane2005} Cochrane, John H. 2005. \emph{Asset Pricing}. Revised edition. Princeton University Press.
\bibitem{donoho2024} Donoho, David. 2024. Data Science at the Singularity. \emph{Harvard Data Science Review}. arXiv:2310.00865.
\bibitem{ff2015} Fama, Eugene F., and Kenneth R. French. 2015. A Five-Factor Asset Pricing Model. \emph{Journal of Financial Economics} 116(1), 1--22.
\bibitem{ctf2025} Hellum, Oliver, Theis Ingerslev Jensen, Bryan T. Kelly, and Lasse Heje Pedersen. 2025. The Power of the Common Task Framework. SSRN 5242901.
\bibitem{jkp2023} Jensen, Theis Ingerslev, Bryan T. Kelly, and Lasse Heje Pedersen. 2023. Is There a Replication Crisis in Finance? \emph{Journal of Finance} 78(5), 2465--2518.
\bibitem{jkpfactors2026} JKP Factors. 2026. Global Factor Data. \url{https://jkpfactors.com/}.
\bibitem{karpathy2026} Karpathy, Andrej. 2026. \emph{autoresearch: AI agents running research on single-GPU nanochat training automatically}. \url{https://github.com/karpathy/autoresearch}.
\bibitem{kx2023} Kelly, Bryan T., and Dacheng Xiu. 2023. Financial Machine Learning. NBER Working Paper 31502.
\bibitem{upsa2023} Kelly, Bryan T., Semyon Malamud, Mohammad Pourmohammadi, and Fabio Trojani. 2023, revised 2025. Universal Portfolio Shrinkage. NBER Working Paper 32004.
\bibitem{markowitz1952} Markowitz, Harry. 1952. Portfolio Selection. \emph{Journal of Finance} 7(1), 77--91.
\bibitem{openai2024o1} OpenAI. 2024. OpenAI o1 System Card. \url{https://openai.com/index/openai-o1-system-card/}.
\bibitem{deepmind2026gemini31} Google DeepMind. 2026. Gemini 3.1 Pro Model Card. \url{https://deepmind.google/models/model-cards/gemini-3-1-pro/}.
\bibitem{xai2025grok41} xAI. 2025. Grok 4.1 Model Card. \url{https://data.x.ai/2025-11-17-grok-4-1-model-card.pdf}.
\bibitem{qwen2025qwen3} Yang, An, et al. 2025. Qwen3 Technical Report. arXiv:2505.09388.
\end{thebibliography}

\end{document}
